% sec1_6.tex — Section 1.6: Generalized Least Squares (GLS)

Assumption~1.4 states that the $n \times n$ matrix of conditional second moments
$\E(\bm{\varepsilon}\bm{\varepsilon}' \mid \vec{X})$ $(= \Var(\bm{\varepsilon} \mid \vec{X}))$
is spherical, that is, proportional to the identity matrix.  Without the assumption,
each element of the $n \times n$ matrix is in general a nonlinear function of $\vec{X}$.
If the error is not (conditionally) homoskedastic, the values of the diagonal elements
of $\E(\bm{\varepsilon}\bm{\varepsilon}' \mid \vec{X})$ are not the same, and if there is
correlation in the error term between observations (the case of serial correlation for
time-series models), the values of the off-diagonal elements are not zero.  For any
given positive scalar $\sigma^2$, define
$\vec{V}(\vec{X}) \equiv \E(\bm{\varepsilon}\bm{\varepsilon}' \mid \vec{X})/\sigma^2$
and assume $\vec{V}(\vec{X})$ is nonsingular and known.  That is,
%
\begin{equation}
  \E(\bm{\varepsilon}\bm{\varepsilon}' \mid \vec{X})
    = \sigma^2 \underset{(n \times n)}{\vec{V}(\vec{X})}, \quad
  \vec{V}(\vec{X}) \text{ nonsingular and known.}
  \label{eq:1.6.1}
\end{equation}

The reason we decompose $\E(\bm{\varepsilon}\bm{\varepsilon}' \mid \vec{X})$ into the
component $\sigma^2$ that is common to all elements of the matrix
$\E(\bm{\varepsilon}\bm{\varepsilon}' \mid \vec{X})$ and the remaining component
$\vec{V}(\vec{X})$ is that we do not need to know the value of $\sigma^2$ for efficient
estimation.  The model that results when Assumption~1.4 is replaced by~\eqref{eq:1.6.1},
which merely assumes that the conditional second moment
$\E(\bm{\varepsilon}\bm{\varepsilon}' \mid \vec{X})$ is nonsingular, is called the
\textbf{generalized regression model}.

%% ─── Consequence of Relaxing Assumption 1.4 ──────────────────────────────────
\subsection*{Consequence of Relaxing Assumption 1.4}

Of the results derived in the previous sections, those that assume Assumption~1.4
are no longer valid for the generalized regression model.  More specifically,
%
\begin{itemize}
  \item The Gauss--Markov Theorem no longer holds for the OLS estimator
    %
    \[
      \vec{b} \equiv (\vec{X}'\vec{X})^{-1}\vec{X}'\vec{y}.
    \]
    %
    The BLUE is some other estimator.

  \item The $t$-ratio is not distributed as the $t$ distribution.  Thus, the $t$-test
    is no longer valid.  The same comments apply to the $F$-test.

  \item However, the OLS estimator \emph{is} still unbiased, because the unbiasedness
    result (Proposition~1.1(a)) does not require Assumption~1.4.
\end{itemize}

%% ─── Efficient Estimation with Known V ───────────────────────────────────────
\subsection*{Efficient Estimation with Known V}

If the value of the matrix function $\vec{V}(\vec{X})$ is known, does there exist a
BLUE for the generalized regression model?  The answer is yes, and the estimator is
called the \textbf{generalized least squares (GLS) estimator}, which we now derive.
The basic idea of the derivation is to transform the generalized regression model,
which consists of Assumptions~1.1--1.3 and~\eqref{eq:1.6.1}, into a model that
satisfies all the assumptions, including Assumption~1.4, of the classical linear
regression model.

For economy of notation, we use $\vec{V}$ for the value $\vec{V}(\vec{X})$.  Since
$\vec{V}$ is by construction symmetric and positive definite, there exists a nonsingular
$n \times n$ matrix $\vec{C}$ such that
%
\begin{equation}
  \vec{V}^{-1} = \vec{C}'\vec{C}.
  \label{eq:1.6.2}
\end{equation}

This decomposition is not unique, with more than one choice for $\vec{C}$, but, as is
clear from the discussion below, the choice of $\vec{C}$ doesn't matter.  Now consider
creating a new regression model by transforming $(\vec{y}, \vec{X}, \bm{\varepsilon})$
by $\vec{C}$ as
%
\begin{equation}
  \tilde{\vec{y}} \equiv \vec{C}\vec{y}, \quad
  \widetilde{\vec{X}} \equiv \vec{C}\vec{X}, \quad
  \tilde{\bm{\varepsilon}} \equiv \vec{C}\bm{\varepsilon}.
  \label{eq:1.6.3}
\end{equation}

Then Assumption~1.1 for $(\vec{y}, \vec{X}, \bm{\varepsilon})$ implies that
$(\tilde{\vec{y}}, \widetilde{\vec{X}}, \tilde{\bm{\varepsilon}})$ too satisfies
linearity:
%
\begin{equation}
  \tilde{\vec{y}} = \widetilde{\vec{X}} \bm{\beta} + \tilde{\bm{\varepsilon}}.
  \label{eq:1.6.4}
\end{equation}

The transformed model satisfies the other assumptions of the classical linear regression
model.  Strict exogeneity is satisfied because
%
\begin{align*}
  \E(\tilde{\bm{\varepsilon}} \mid \widetilde{\vec{X}})
    &= \E(\tilde{\bm{\varepsilon}} \mid \vec{X}) \\
    &\qquad \text{(since $\vec{C}$ is nonsingular, $\vec{X}$ and
      $\widetilde{\vec{X}}$ contain the same information)} \\
    &= \E(\vec{C}\bm{\varepsilon} \mid \vec{X}) \\
    &= \vec{C}\,\E(\bm{\varepsilon} \mid \vec{X})
      \quad \text{(by the linearity of conditional expectations)} \\
    &= \vec{0}
      \quad \text{(since $\E(\bm{\varepsilon} \mid \vec{X}) = \vec{0}$
        by Assumption~1.2).}
\end{align*}

Because $\vec{V}$ is positive definite, the no-multicollinearity assumption is also
satisfied (see a review question below for a proof).  Assumption~1.4 is satisfied for
the transformed model because
%
\begin{align*}
  \E(\tilde{\bm{\varepsilon}}\tilde{\bm{\varepsilon}}' \mid \widetilde{\vec{X}})
    &= \E(\tilde{\bm{\varepsilon}}\tilde{\bm{\varepsilon}}' \mid \vec{X})
      \quad \text{(since $\widetilde{\vec{X}}$ and $\vec{X}$ contain the same
        information)} \\
    &= \vec{C}\,\E(\bm{\varepsilon}\bm{\varepsilon}' \mid \vec{X})\,\vec{C}'
      \quad \text{(since
        $\tilde{\bm{\varepsilon}}\tilde{\bm{\varepsilon}}'
          = \vec{C}\bm{\varepsilon}\bm{\varepsilon}'\vec{C}'$)} \\
    &= \vec{C} \cdot \sigma^2 \cdot \vec{V}\vec{C}'
      \quad \text{(by~\eqref{eq:1.6.1})} \\
    &= \sigma^2 \vec{C}\vec{V}\vec{C}' \\
    &= \sigma^2 \vec{I}_n
      \quad \text{(since $(\vec{C}')^{-1}\vec{V}^{-1}\vec{C}^{-1} = \vec{I}_n$
        or $\vec{C}\vec{V}\vec{C}' = \vec{I}_n$ by~\eqref{eq:1.6.2}).}
\end{align*}

So indeed the variance of the transformed error vector $\tilde{\bm{\varepsilon}}$ is
spherical.  Finally, $\tilde{\bm{\varepsilon}} \mid \widetilde{\vec{X}}$ is normal
because the distribution of $\tilde{\bm{\varepsilon}} \mid \widetilde{\vec{X}}$ is
the same as $\tilde{\bm{\varepsilon}} \mid \vec{X}$ and $\tilde{\bm{\varepsilon}}$
is a linear transformation of $\bm{\varepsilon}$.  This completes the verification of
Assumptions~1.1--1.5 for the transformed model.

The Gauss--Markov Theorem for the transformed model implies that the BLUE of
$\bm{\beta}$ for the generalized regression model is the OLS estimator applied
to~\eqref{eq:1.6.4}:
%
\begin{align}
  \widehat{\bm{\beta}}_{\text{GLS}}
    &= (\widetilde{\vec{X}}'\widetilde{\vec{X}})^{-1}
       \widetilde{\vec{X}}'\tilde{\vec{y}} \notag \\
    &= [(\vec{C}\vec{X})'(\vec{C}\vec{X})]^{-1}(\vec{C}\vec{X})'\vec{C}\vec{y}
       \notag \\
    &= (\vec{X}'\vec{C}'\vec{C}\vec{X})^{-1}(\vec{X}'\vec{C}'\vec{C}\vec{y})
       \notag \\
    &= (\vec{X}'\vec{V}^{-1}\vec{X})^{-1}\vec{X}'\vec{V}^{-1}\vec{y}
       \quad \text{(by~\eqref{eq:1.6.2}).}
  \label{eq:1.6.5}
\end{align}

This is the GLS estimator.  Its conditional variance is
%
\begin{align}
  \Var(\widehat{\bm{\beta}}_{\text{GLS}} \mid \vec{X})
    &= (\vec{X}'\vec{V}^{-1}\vec{X})^{-1}\vec{X}'\vec{V}^{-1}
       \Var(\vec{y} \mid \vec{X})\vec{V}^{-1}\vec{X}
       (\vec{X}'\vec{V}^{-1}\vec{X})^{-1} \notag \\
    &= (\vec{X}'\vec{V}^{-1}\vec{X})^{-1}\vec{X}'\vec{V}^{-1}
       (\sigma^2 \vec{V})\vec{V}^{-1}\vec{X}
       (\vec{X}'\vec{V}^{-1}\vec{X})^{-1} \notag \\
    &\qquad \text{(since $\Var(\vec{y} \mid \vec{X}) = \Var(\bm{\varepsilon}
      \mid \vec{X})$)} \notag \\
    &= \sigma^2 \cdot (\vec{X}'\vec{V}^{-1}\vec{X})^{-1}.
  \label{eq:1.6.6}
\end{align}

Since replacing $\vec{V}$ by $\sigma^2 \cdot \vec{V}$ $(= \Var(\bm{\varepsilon}
\mid \vec{X}))$ in~\eqref{eq:1.6.5} does not change the numerical value, the GLS
estimator can also be written as
%
\[
  \widehat{\bm{\beta}}_{\text{GLS}}
    = \bigl[\vec{X}'\,\Var(\bm{\varepsilon} \mid \vec{X})^{-1}\,\vec{X}\bigr]^{-1}
      \vec{X}'\,\Var(\bm{\varepsilon} \mid \vec{X})^{-1}\,\vec{y}.
\]

As noted above, the OLS estimator $(\vec{X}'\vec{X})^{-1}\vec{X}'\vec{y}$ too is
unbiased without Assumption~1.4, but nevertheless the GLS estimator should be preferred
(provided $\vec{V}$ is known) because the latter is more efficient in that the variance
is smaller in the matrix sense.  The gain in efficiency is achieved by exploiting the
heteroskedasticity and correlation between observations in the error term, which,
operationally, is to insert the inverse of (a matrix proportional to)
$\Var(\bm{\varepsilon} \mid \vec{X})$ in the OLS formula, as in~\eqref{eq:1.6.5}.
The discussion so far can be summarized as

\begin{proposition}[finite-sample properties of GLS]
\label{prop:1.7}
\mbox{}
\begin{enumerate}[label=(\alph*)]
  \item \emph{(unbiasedness)}\enspace
    Under Assumptions~1.1--1.3,
    $\E(\widehat{\bm{\beta}}_{\textup{GLS}} \mid \vec{X}) = \bm{\beta}$.

  \item \emph{(expression for the variance)}\enspace
    Under Assumptions~1.1--1.3 and the assumption~\eqref{eq:1.6.1} that the
    conditional second moment is proportional to $\vec{V}(\vec{X})$,
    %
    \[
      \Var(\widehat{\bm{\beta}}_{\textup{GLS}} \mid \vec{X})
        = \sigma^2 \cdot (\vec{X}'\vec{V}(\vec{X})^{-1}\vec{X})^{-1}.
    \]

  \item \emph{(efficiency of GLS)}\enspace
    Under the same set of assumptions as in~(b), the GLS estimator is efficient
    in that the conditional variance of any unbiased estimator that is linear in
    $\vec{y}$ is greater than or equal to
    $\Var(\widehat{\bm{\beta}}_{\textup{GLS}} \mid \vec{X})$ in the matrix sense.
\end{enumerate}
\end{proposition}

%% ─── A Special Case: Weighted Least Squares (WLS) ────────────────────────────
\subsection*{A Special Case: Weighted Least Squares (WLS)}

The idea of adjusting for the error variance matrix becomes more transparent when
there is no correlation in the error term between observations so that the matrix
$\vec{V}$ is diagonal.  Let $v_i(\vec{X})$ be the $i$-th diagonal element of
$\vec{V}(\vec{X})$.  So
%
\[
  \E(\varepsilon_i^2 \mid \vec{X})
    \;(= \Var(\varepsilon_i \mid \vec{X}))
    = \sigma^2 \cdot v_i(\vec{X}).
\]

It is easy to see that $\vec{C}$ is also diagonal, with the square root of
$1/v_i(\vec{X})$ in the $i$-th diagonal.  Thus,
$(\tilde{\vec{y}}, \widetilde{\vec{X}})$ is given by
%
\[
  \tilde{y}_i = \frac{y_i}{\sqrt{v_i(\vec{X})}}, \quad
  \tilde{\vec{x}}_i = \frac{\vec{x}_i}{\sqrt{v_i(\vec{X})}}
  \quad (i = 1, 2, \ldots, n).
\]

Therefore, efficient estimation under a known form of heteroskedasticity is first to
weight each observation by the reciprocal of the square root of the variance
$v_i(\vec{X})$ and then apply OLS\@.  This is called the \textbf{weighted regression}
(or the \textbf{weighted least squares (WLS)}).

An important further special case is the case of a random sample where
$\{y_i, \vec{x}_i\}$ is i.i.d.\ across $i$.  As was noted in Section~1.1, the error
is unconditionally homoskedastic (i.e., $\E(\varepsilon_i^2)$ does not depend on $i$),
but still GLS can be used to increase efficiency because the error can be conditionally
heteroskedastic.  The conditional second moment
$\E(\varepsilon_i^2 \mid \vec{X})$ for the case of random samples depends only on
$\vec{x}_i$, and the functional form of $\E(\varepsilon_i^2 \mid \vec{x}_i)$ is the
same across $i$.  Thus
%
\begin{equation}
  v_i(\vec{X}) = v(\vec{x}_i) \quad \text{for random samples.}
  \label{eq:1.6.7}
\end{equation}

So the knowledge of $\vec{V}(\cdot)$ comes down to a single function of $K$ variables,
$v(\cdot)$.

%% ─── Limiting Nature of GLS ─────────────────────────────────────────────────
\subsection*{Limiting Nature of GLS}

All these sanguine conclusions about the finite-sample properties of GLS rest on the
assumption that the regressors in the generalized regression model are strictly
exogenous ($\E(\tilde{\bm{\varepsilon}} \mid \widetilde{\vec{X}}) = \vec{0}$).  This
fact limits the usefulness of the GLS procedure.  Suppose, as is often the case with
time-series models, that the regressors are not strictly exogenous and the error is
serially correlated.  So neither OLS nor GLS has those good finite-sample properties
such as unbiasedness.  Nevertheless, as will be shown in the next chapter, the OLS
estimator, which ignores serial correlation in the error, will have some good
large-sample properties (such as ``consistency'' and ``asymptotic normality''),
provided that the regressors are ``predetermined'' (which is weaker than strict
exogeneity).  The GLS estimator, in contrast, does not have that redeeming feature.
That is, if the error is not strictly exogenous but is merely predetermined, the GLS
procedure to correct for serial correlation can make the estimator inconsistent (see
Section~6.7).  A procedure for explicitly taking serial correlation into account while
maintaining consistency will be presented in Chapter~6.

If it is not appropriate for correcting for serial correlation, the GLS procedure can
still be used to correct for heteroskedasticity when the error is not serially
correlated with diagonal $\vec{V}(\vec{X})$, in the form of WLS\@.  But that is
provided that the matrix function $\vec{V}(\vec{X})$ is known.  Very rarely do we have
\emph{a priori} information specifying the values of the diagonal elements of
$\vec{V}(\vec{X})$, which is necessary to weight observations.  In the case of a
random sample where serial correlation is guaranteed not to arise, the knowledge of
$\vec{V}(\vec{X})$ boils down to a single function of $K$ variables, $v(\vec{x}_i)$,
as we have just seen, but even for this case the knowledge of such a function is
unavailable in most applications.

If we do not know the function $\vec{V}(\vec{X})$, we can estimate its functional form
from the sample.  This approach is called the \textbf{Feasible Generalized Least Squares
(FGLS)}.  But if the function $\vec{V}(\vec{X})$ is estimated from the sample, its
value $\vec{V}$ becomes a random variable, which affects the distribution of the GLS
estimator.  Very little is known about the finite-sample properties of the FGLS
estimator.  We will cover the large-sample properties of the FGLS estimator in the
context of heteroskedasticity correction in the next chapter.

Before closing, one positive side of GLS should be noted: most linear estimation
techniques---including the 2SLS, 3SLS, and the random effects estimators to be
introduced later---can be expressed as a GLS estimator, with some liberal definition
of data matrices.  However, those estimators and OLS can also be interpreted as a GMM
(generalized method of moments) estimator, and the GMM interpretation is more useful
for developing large-sample results.
